
\documentclass{article}
\usepackage{colortbl}
\usepackage{makecell}
\usepackage{multirow}
\usepackage{supertabular}

\begin{document}

\newcounter{utterance}

\centering \large Interaction Transcript for game `function\_detective', experiment `logic\_passive\_labeled\_pairs\_oneshot', episode 191 with mistral{-}large{-}3{-}azure.
\vspace{24pt}

{ \footnotesize  \setcounter{utterance}{1}
\setlength{\tabcolsep}{0pt}
\begin{supertabular}{c@{$\;$}|p{.15\linewidth}@{}p{.15\linewidth}p{.15\linewidth}p{.15\linewidth}p{.15\linewidth}p{.15\linewidth}}
    \# & \multicolumn{2}{c}{Player} && \multicolumn{2}{c}{Game Master} \\
    \hline

    \theutterance \stepcounter{utterance}  
    & & & \multicolumn{4}{p{0.6\linewidth}}{
        \cellcolor[rgb]{0.9,0.9,0.9}{
            \makecell[{{p{\linewidth}}}]{
                \texttt{\tiny{[P1$\langle$GM]}}
                \texttt{I have a hidden Python function for you to figure out.} \\
\\ 
\texttt{Your job is to infer what it does and then submit a correct Python implementation.} \\
\\ 
\texttt{Function signature:} \\
\texttt{(a: bool, b: bool) {-}> bool} \\
\\ 
\texttt{You have 1 turns total.} \\
\\ 
\texttt{I have already given you a fixed list of candidate input{-}output pairs with true or false labels. Solve using those statements only.} \\
\\ 
\texttt{You must make your best SOLVE attempt no later than the second{-}to{-}last turn.} \\
\texttt{Do not wait until the very end.} \\
\\ 
\texttt{You can use these actions:} \\
\texttt{{-} JSON format described at the end of the prompt} \\
\\ 
\\ 
\texttt{Do not add import statements in your guess.} \\
\texttt{Anything you might need is already available, so just write the function body in hypothesis and return the answer.} \\
\\ 
\texttt{Preloaded examples:} \\
\texttt{1. (((True, False), False), True)} \\
\texttt{2. (((True, True), True), False)} \\
\texttt{3. (((False, True), True), True)} \\
\texttt{4. (((False, False), True), False)} \\
\texttt{5. (((True, False), False), True)} \\
\texttt{6. (((True, True), True), False)} \\
\texttt{7. (((False, True), True), True)} \\
\texttt{8. (((False, False), True), False)} \\
\texttt{9. (((True, False), False), True)} \\
\texttt{10. (((True, True), True), False)} \\
\texttt{11. (((False, True), True), True)} \\
\texttt{12. (((False, False), True), False)} \\
\texttt{13. (((True, False), False), True)} \\
\texttt{14. (((True, True), True), False)} \\
\texttt{15. (((False, True), True), True)} \\
\texttt{16. (((False, False), True), False)} \\
\texttt{17. (((True, False), False), True)} \\
\texttt{18. (((True, True), True), False)} \\
\texttt{19. (((False, True), True), True)} \\
\texttt{20. (((False, False), True), False)} \\
\texttt{21. (((True, False), False), True)} \\
\texttt{22. (((True, True), True), False)} \\
\texttt{23. (((False, True), True), True)} \\
\texttt{24. (((False, False), True), False)} \\
\texttt{25. (((True, False), False), True)} \\
\texttt{26. (((True, True), True), False)} \\
\texttt{27. (((False, True), True), True)} \\
\texttt{28. (((False, False), True), False)} \\
\texttt{29. (((True, False), False), True)} \\
\texttt{30. (((True, True), True), False)} \\
\\ 
\texttt{Scoring:} \\
\texttt{{-} correctness is required for any non{-}zero score} \\
\texttt{{-} in interactive modes, solving with fewer information requests is better} \\
\\ 
\texttt{Input formatting reminders:} \\
\texttt{{-} strings must be quoted} \\
\texttt{{-} booleans must be True/False} \\
\texttt{{-} use commas between multiple values} \\
\texttt{{-} do not write any extra action tags in the same message} \\
\\ 
\texttt{Format example:} \\
\texttt{You write:} \\
\texttt{\{'current\_mode': <probing/confirming/don't know; choose whichever best describes your current state in the task>, 'input\_rationale': <brief description of why you chose this hypothesis>, 'hypothesis': <either the exact string "unknown" if you have no guess yet OR your current best guess as runnable Python matching the given signature, with a one{-}line docstring stating what it does. Never write placeholder code, if you cannot commit to a guess, use "unknown". Match the given signature.} \\
\texttt{e.g.} \\
\texttt{def f(x):} \\
\texttt{	"""Returns the absolute value of x."""} \\
\texttt{	return x if x >= 0 else {-}x} \\
\\ 
\texttt{Use "unknown" if you have no guess.>, 'confidence\_score': <score between 1{-}100. 100 means highest confidence in your hypothesis>\}} \\
\\ 
\texttt{Output only one valid json per turn.} \\
            }
        }
    }
     \\ \\

    \theutterance \stepcounter{utterance}  
    & \multicolumn{4}{p{0.6\linewidth}}{
        \cellcolor[rgb]{0.561,0.737,0.561}{
            \makecell[{{p{\linewidth}}}]{
                \texttt{\tiny{[P1$\rangle$GM]}}
                \texttt{```json} \\
\texttt{\{} \\
\texttt{  "current\_mode": "confirming",} \\
\texttt{  "input\_rationale": "Analyzing the provided input{-}output pairs, the function seems to return True only when the first argument (a) is True and the second argument (b) is False, or when the first argument (a) is False and the second argument (b) is True. This matches the behavior of the logical XOR (exclusive OR) operation, but with a twist: the output is negated when the second input to the function (the 'b' in the signature) is True. However, upon closer inspection, the function appears to return True if 'a' is not equal to the second input (the 'b' in the function call, not the 'b' in the signature). Wait, no—the examples show that the function's output depends on the relationship between 'a' and the first element of the tuple input, and 'b' and the second element of the tuple input. Actually, the function seems to return True if 'a' is not equal to the first element of the tuple and 'b' is not equal to the second element of the tuple. But that doesn't fit all cases. Alternatively, the function might be returning True if 'a' is not equal to the second input (the 'b' in the function call). Given the examples, the function returns True when 'a' is True and the second input (the 'b'} \\
            }
        }
    }
    & & \\ \\

    \theutterance \stepcounter{utterance}  
    & & & \multicolumn{2}{p{0.3\linewidth}}{
        \cellcolor[rgb]{0.9,0.9,0.9}{
            \makecell[{{p{\linewidth}}}]{
                \texttt{\tiny{[GM$|$GM]}}
                \texttt{Game Verdict: LOSS (turn limit reached)} \\
            }
        }
    }
    & & \\ \\

\end{supertabular}
}

\end{document}
